\chapter{Metodyka}\label{ch:metodyka}

\TemplateNote{Rozdział pokazuje wzory, pseudokod (środowisko \code{algorithm}) oraz zwykłą tabelę. Tabele mają podpis
NAD tabelą, jeden rozmiar czcionki (\latexcmd{tablefont}) i styl \code{booktabs}. Nie skaluj tabel poleceniem
\latexcmd{resizebox} -- do rozciągania służy środowisko \code{tabularx} i kolumny \code{Y} oraz \code{Z}.}

\section{Założenia}\label{sec:zalozenia}

Realizację zadania oparto na następujących założeniach:
\begin{enumerate}
    \item wszystkie czynności wykonywane są w powtarzalny sposób i opisane z podaniem użytych narzędzi,
    \item wyniki pośrednie są zapisywane, a ich poprawność weryfikowana (np. skrótami kryptograficznymi),
    \item każdy rysunek, tabela i listing mają podpis oraz etykietę i są przywoływane w tekście.
\end{enumerate}

\section{Opis metody}\label{sec:opis-metody}

Wzory wyróżnione numeruje się środowiskiem \code{equation}, a wzory wieloliniowe -- środowiskiem \code{align}.
Przykładowo, średnia arytmetyczna $n$ pomiarów $x_1, \dots, x_n$ wyraża się wzorem~\eqref{eq:srednia}, a jej
odchylenie standardowe wzorem~\eqref{eq:odchylenie}:
\begin{align}
    \bar{x} &= \frac{1}{n} \sum_{i=1}^{n} x_i , \label{eq:srednia} \\
    s       &= \sqrt{\frac{1}{n-1} \sum_{i=1}^{n} \left( x_i - \bar{x} \right)^2} . \label{eq:odchylenie}
\end{align}

Wzory krótkie, występujące w tekście, zapisuje się w znakach dolara, np. $\RR$, $\NN$ lub $a \IMP b$
(skróty zdefiniowane w \file{preamble/macros.tex}).

\subsection{Algorytm wyszukiwania}\label{sec:algorytm}

Pseudokod przedstawiono w algorytmie~\ref{alg:binarne}. Środowisko automatycznie numeruje algorytmy w obrębie
rozdziału i dopisuje je do spisu algorytmów.

\begin{algorithm}[htbp]
    \caption{Wyszukiwanie binarne w tablicy posortowanej rosnąco}\label{alg:binarne}
    \Input{posortowana tablica $A[1..n]$, szukana wartość $x$}
    \Output{indeks elementu równego $x$ albo $0$, gdy takiego elementu nie ma}
    $l \gets 1$; \quad $p \gets n$\;
    \While{$l \leq p$}{
        $s \gets \lfloor (l + p) / 2 \rfloor$\;
        \eIf{$A[s] = x$}{
            \Return $s$\;
        }{
            \eIf{$A[s] < x$}{
                $l \gets s + 1$\;
            }{
                $p \gets s - 1$\;
            }
        }
    }
    \Return $0$\;
\end{algorithm}

\section{Zestawienie parametrów}\label{sec:parametry}

Tabela~\ref{tab:parametry} przedstawia typową tabelę z kolumną rozciągliwą (\code{Y}) i kolumnami o stałej
szerokości. Komórki nagłówka oznacza się makrem \latexcmd{hcell}.

\begin{table}[H]
    \centering
    \caption{Przykładowe zestawienie parametrów przyjętych w badaniu}\label{tab:parametry}
    \tablefont
    \begin{tabularx}{\linewidth}{l Y l}
        \toprule
        \hcell{Parametr} & \hcell{Opis} & \hcell{Wartość} \\
        \midrule
        $n$      & liczba pomiarów w pojedynczej serii       & 100 \\
        $\alpha$ & poziom istotności przyjęty w testach      & 0,05 \\
        $T$      & maksymalny czas trwania pojedynczej próby & 30\,s \\
        \bottomrule
    \end{tabularx}
\end{table}
